\documentclass[11pt]{article}
\usepackage{amsmath}
\usepackage{amssymb}
\usepackage{tabularx}
\usepackage{booktabs}
\usepackage{graphicx}
\usepackage{multirow}
\usepackage{array}
\usepackage[hidelinks]{hyperref}
\usepackage[T1]{fontenc}
\usepackage[utf8]{inputenc}
\usepackage{microtype}

\title{Appendix: Detailed Implementation, Experiments, and Results}
\author{Daniya Niazi}
\date{}

\begin{document}

\maketitle
\tableofcontents

\newpage

\section{Dataset Details and Characteristics}

This section provides comprehensive information about the four medical imaging datasets used in our certification study, including modality specifications, dataset sizes, class distributions, and representative sample images.

\subsection{ChestX-ray14}

ChestX-ray14 consists of frontal chest radiographs (X-ray, transmission imaging) at 1024×1024 pixels (original, downsampled to 224×224), grayscale, with 16-bit depth preserving full radiographic detail. The dataset contains 112,120 total images with an 80/10/10 train/val/test split, yielding 100 test images for Experiment 1 and 5 per architecture for Experiment 2.

Classification involves 14 thoracic pathology classes: atelectasis, cardiomegaly, effusion, infiltration, mass, nodule, pneumonia, pneumothorax, consolidation, edema, emphysema, fibrosis, pleural thickening, and hernia. The training set shows class imbalance typical of clinical datasets: 67\% normal findings, 9\% pneumonia, 7\% effusion, 4\% cardiomegaly, 3\% infiltration, and 10\% other conditions. Multi-label annotations reflect that single images often contain multiple pathologies.

X-rays represent the most commonly ordered medical imaging test globally. Learning to certify explanations on chest X-rays directly impacts high-volume, safety-critical deployment. Representative samples span normal anatomy, pneumonia with infiltrates, cardiomegaly with enlarged cardiac silhouettes, and pleural effusion with fluid collections.

\subsection{Brain MRI Tumor Dataset}

Brain MRI consists of T1-weighted magnetic resonance images with variable resolution (150–512 pixels, resized to 224×224), grayscale, slice thickness 1–2 mm, and 90° flip angle for T1 contrast. The dataset contains 3,064 multi-slice volumes with 80/10/10 split, providing 100 test images for Experiment 1.

Four-class classification distinguishes glioma (most common primary tumor), meningioma (usually benign), pituitary tumor (small endocrine lesions), and normal brain. Training class distribution is well-balanced: 34\% glioma, 25\% meningioma, 25\% normal, 16\% pituitary. Brain tumors are life-threatening; accurate diagnosis and localization are critical for surgical planning. Our certification study validates whether AI explanations reliably highlight tumor regions. Representative samples include normal brain anatomy, glioma with infiltrative appearance, well-defined meningioma masses, and small pituitary lesions.

\subsection{APTOS 2019 (Diabetic Retinopathy)}

APTOS 2019 consists of fundus photographs (color reflection imaging) at 1024×1024 or 2048×2048 pixels (resized to 224×224), RGB color space, 45° field of view, with mydriatic imaging. The dataset contains 3,662 images with 80/10/10 split, yielding 100 test images for experiments.

Classification uses five-class ordinal severity grading: Grade 0 (no DR, normal retina), Grade 1 (mild: microaneurysms and hemorrhages), Grade 2 (moderate: more microaneurysms, cotton-wool spots, hard exudates), Grade 3 (severe: numerous hemorrhages, venous beading, intraretinal microvascular abnormalities), and Grade 4 (proliferative: neovascularization, vitreous hemorrhage, tractional detachment). Training distribution reflects clinical prevalence: 49\% no DR, 20\% moderate, 13\% mild, 11\% severe, 7\% proliferative.

Diabetic retinopathy is the leading cause of blindness in working-age adults. Early detection requires identifying subtle vascular changes; color information is essential for certifying color-based explanations in clinical deployment. Representative samples span normal retina with optic disc, mild DR with microaneurysms, moderate DR with hemorrhages and exudates, and proliferative DR with neovascularization.

\subsection{ISIC Skin Lesion Dataset}

ISIC consists of dermoscopic images (magnified color photography) at 600×450 to 6000×4000 pixels (resized to 224×224), RGB color space, 10×–40× magnification, with polarized or contact lighting. The dataset contains 13,900 images with 80/10/10 split.

Seven-class classification distinguishes melanoma (malignant, high mortality), melanocytic nevus (common benign mole), basal cell carcinoma (common, low mortality), actinic keratosis (precancerous), benign keratosis (seborrheic), dermatofibroma (benign fibrous), and vascular lesion (hemangioma, angioma). Training distribution shows severe class imbalance: 60\% melanocytic nevi, 10\% melanoma, 10\% benign keratosis, 5\% basal cell carcinoma, 3\% actinic keratosis, 3\% vascular lesion, 1\% dermatofibroma.

Melanoma is the deadliest skin cancer; early detection dramatically improves survival. However, dermoscopic images present high visual complexity with color variation, texture patterns, and artifacts (hair, rulers). This dataset represents the most challenging for certification. Key challenges include hair artifacts, ruler/scale presence, skin markings, high color variation, and class imbalance. Representative samples span benign nevi with regular patterns, melanoma with irregular borders and color variation, basal cell carcinoma with pearly appearance, and images with hair artifacts.

\newpage

\section{Model Architecture Selection and Training Results}

This section describes the six CNN architectures trained and tuned on each medical imaging dataset, along with their validation and test accuracies.

\subsection{Architecture Specifications}

Six CNN architectures were trained on each medical imaging dataset. ResNet-18, ResNet-34, and ResNet-50 employ skip connections enabling gradient flow, with 11.7M, 21.8M, and 25.6M parameters respectively and 4 residual blocks. EfficientNet-B0 and B1 use compound scaling (depth, width, resolution jointly) with 5.3M and 7.8M parameters. DenseNet-121 uses dense connections with feature reuse (7.9M parameters). All models were pre-trained on ImageNet-1k, then fine-tuned on medical datasets using Adam optimizer (learning rate 1e-4), batch size 32, 30 epochs with early stopping (patience=5), and standard augmentation (rotation ±15°, horizontal flip, color jitter for RGB only).

\subsection{Training Results by Dataset}

\subsubsection{ChestX-ray14}

EfficientNet-B1 achieved highest test accuracy (90.5\%) on ChestX-ray14, followed by DenseNet-121 (89.8\%), ResNet-50 (89.1\%), EfficientNet-B0 (89.4\%), ResNet-34 (88.9\%), and ResNet-18 (88.4\%). Validation-test gaps remained below 1\%, indicating no significant overfitting. The multi-label nature of the task (multiple pathologies per image) challenges all models. ResNet-18 was selected for comprehensive Experiment 1 due to reasonable accuracy and faster certification.

\subsubsection{Brain MRI Tumor}

Brain MRI achieved highest overall accuracies, with EfficientNet-B1 at 92.7\%, followed by DenseNet-121 (92.3\%), ResNet-50 (91.7\%), EfficientNet-B0 (92.0\%), ResNet-34 (91.4\%), and ResNet-18 (90.9\%). The four-class problem is more balanced than ChestX-ray, and the grayscale nature may contribute to cleaner discriminative features. ResNet-18 remains a strong performer despite architectural simplicity.

\subsubsection{APTOS 2019 (Diabetic Retinopathy)}

APTOS achieved lower accuracies than grayscale datasets, with EfficientNet-B1 at 88.7\%, DenseNet-121 (88.1\%), EfficientNet-B0 (87.9\%), ResNet-50 (87.6\%), ResNet-34 (87.3\%), and ResNet-18 (86.8\%). Color complexity and fine vascular structures make this task more challenging. High Kappa scores (0.82–0.84) indicate strong compatibility with inter-rater clinical agreement.

\subsubsection{ISIC Skin Lesion}

ISIC achieved lowest overall accuracies due to class imbalance and artifact complexity. EfficientNet-B1 reached 86.9\%, DenseNet-121 (86.2\%), EfficientNet-B0 (85.9\%), ResNet-50 (85.5\%), ResNet-34 (85.1\%), and ResNet-18 (84.4\%). High color variation, fine texture details, and the 60\%-10\% melanoma-nevus imbalance complicate learning. Despite lower accuracy, these models provide baseline comparisons for certification analysis.

\subsection{Cross-Dataset Model Performance Summary}

Average test accuracy across all datasets reveals systematic trends: EfficientNet-B1 achieved highest average (89.7\%), followed by DenseNet-121 (89.1\%), ResNet-50 (88.5\%), EfficientNet-B0 (88.8\%), ResNet-34 (88.2\%), and ResNet-18 (87.6\%). All architectures show consistent modality ranking: grayscale ChestX-ray highest, followed by Brain MRI, APTOS fundus, then ISIC skin with color artifacts and imbalance. EfficientNet-B1's accuracy advantage masks complex attribution behavior that emerges in certification analysis.

\begin{table}[h!]
\centering
\caption{Model test accuracy snapshot (best and runner-up per dataset)}
\begin{tabularx}{\linewidth}{|X|X|X|}
\hline
	extbf{Dataset} & \textbf{Best model (Test Acc)} & \textbf{Runner-up} \\
\hline
ChestX-ray14 & EfficientNet-B1 (90.5\%) & DenseNet-121 (89.8\%) \\
\hline
Brain MRI & EfficientNet-B1 (92.7\%) & DenseNet-121 (92.3\%) \\
\hline
APTOS DR & EfficientNet-B1 (88.7\%) & DenseNet-121 (88.1\%) \\
\hline
ISIC Skin & EfficientNet-B1 (86.9\%) & DenseNet-121 (86.2\%) \\
\hline
\end{tabularx}
\label{tab:acc_summary}
\end{table}

\newpage

\section{Attribution Methods: Settings and Implementation}

This section details the configuration and implementation of each attribution method, along with visual results on representative samples.

\subsection{Integrated Gradients}

Integrated Gradients computes path integration of gradients from a black-image baseline to the input: $IG_i(x) = (x_i - x'_i) \int_0^1 \frac{\partial f(\alpha x + (1-\alpha) x')}{\partial x_i} d\alpha$, where $x'$ is the baseline. Implementation uses 50 integration steps with Riemann sum approximation and softmax normalization to [0,1]. Computational cost is 3–5 seconds per image (50 forward passes required). Smooth, diffuse attributions naturally smooth via path integration, identify global disease-relevant regions, and prove robust to small perturbations—explaining strong certification performance. Integrated Gradients achieves 52–68\% certification across datasets (highest among methods).

\subsection{GradCAM (Gradient-weighted Class Activation Map)}

GradCAM combines gradients with convolutional layer activations via $GradCAM_c = ReLU\left(\sum_k \alpha_k^c A^k\right)$, where $\alpha_k^c = \frac{1}{Z} \sum_{i,j} \frac{\partial y^c}{\partial A_{ij}^k}$ measures gradient of class score relative to activation. Implementation targets layer4 (final residual block before pooling) with bilinear upsampling to 224×224 and ReLU to eliminate negative values. Computational cost is 2–4 seconds per image. Coarse spatial resolution (14×14 upsampled) acts as implicit smoothing, producing strong focus on discriminative features. GradCAM achieves 38–55\% certification, moderate performance reflecting coarse resolution benefits but some loss of localization precision.

\subsection{Layer-wise Relevance Propagation (LRP)}

LRP distributes prediction score backward through layers via $R_j^{(l)} = \sum_k \frac{a_j^{(l)} w_{jk}^{(l,l+1)}}{\sum_j a_j^{(l)} w_{jk}^{(l,l+1)} + \epsilon} R_k^{(l+1)}$, ensuring relevance is neither created nor destroyed. Implementation uses $\epsilon$-rule ($\epsilon=0.01$) for fully connected layers and $z^+$-rule for convolutional layers. Computational cost is 5–8 seconds per image. Attributions often appear sharper than Integrated Gradients, highlighting precise feature locations, but show high variability—excellent on some images, poor on others. This sensitivity to ReLU discontinuities explains the 32–52\% certification range and makes LRP unreliable without explicit robustness constraints.

\subsection{RISE (Randomized Input Sampling for Explanation)}

RISE generates random binary masks ($N_{masks}$ total), observes prediction changes, and computes weighted averages. The algorithm: (1) generates $N_{masks}$ random binary masks (resolution $s \times s$), (2) upsamples masks bilinearly and applies to image, (3) predicts on masked images, (4) weights pixels by average confidence drop across masks. Implementation reduces masks from typical 8000 to 500 for computational feasibility, uses 7×7 mask resolution with bilinear upsampling, and measures maximum neuron value across output (class-agnostic). Computational cost in certification: 500 masks × 100 noise samples = 45–60 seconds per image. Coarse, blocky attributions reflect mask resolution limitations; high random sampling variance amplifies under certification. Reduced mask count significantly impacts quality, explaining poor 20–28\% certification rates.

\subsection{Occlusion}

Occlusion systematically blocks image patches and measures prediction changes. Algorithm: (1) define patch size $P$ and stride $S$, (2) for each patch position: zero out patch, predict, measure confidence drop, (3) assign patch importance based on confidence change. Implementation uses 16×16 patches, 8-pixel stride (50\% overlap), yielding 28×28 grid for 224×224 images, measuring softmax score drop for predicted class. Computational cost: ~784 forward passes (28×28 + original) = 80–120 seconds per image. Discrete, patch-aligned attributions subject to patch alignment sensitivity—small perturbations shift patches, changing attribution maps. High computational cost and poor certification (12–22\%) reflect both expensive evaluation and fundamental brittleness of patch-based approaches.

\newpage

\section{Experiment 1 Results: Comprehensive Certification on ResNet-18}

This section presents detailed results from Experiment 1, certifying attributions on 100 test images per dataset using ResNet-18 architecture across five attribution methods and varying $\sigma$ and $N$ parameters.

\subsection{Certification Rates by Dataset and Method}

\subsubsection{ChestX-ray14 Certification Results}

ChestX-ray14 demonstrates strong certification rates across methods. Integrated Gradients achieves 68\% certification at moderate noise ($\sigma=0.5$, $N=100$), with GradCAM at 55\%, LRP at 52\%, RISE at 28\%, and Occlusion at 22\%. Grayscale imaging provides certification advantage through homogeneous lung fields enabling smooth attributions. Pathologies (infiltrates, masses) create strong, certifiable signals. Performance improves with increased noise standard deviation ($\sigma$), trading localization precision for certification guarantees.

\subsubsection{Brain MRI Tumor Certification Results}

Brain MRI achieves Integrated Gradients certification of 62\% at $\sigma=0.5$, with GradCAM at 48\%, LRP at 45\%, RISE at 25\%, and Occlusion at 18\%. Well-defined tumor masses certify better than diffuse changes like edema. Grayscale again provides certification advantage. Subtle texture differences in surrounding edema challenge certification. High clinical relevance confirmed by 68\% IoU overlap with radiologist tumor masks.

\subsubsection{APTOS 2019 (Fundus) Certification Results}

APTOS fundus images show 58\% certification for Integrated Gradients at $\sigma=0.5$, with GradCAM at 42\%, LRP at 38\%, RISE at 22\%, and Occlusion at 15\%. Color channels reduce certification rates by 10–15\% versus grayscale datasets. Fine vascular structures (microaneurysms) prove highly sensitive to smoothing. Certified pixels localize to major hemorrhages and exudates when present, but miss subtle microvascular changes critical for clinical staging.

\subsubsection{ISIC Skin Lesion Certification Results}

ISIC skin lesions present the most challenging certification scenario: 52\% for Integrated Gradients at $\sigma=0.5$, GradCAM 38\%, LRP 32\%, RISE 20\%, Occlusion 12\%. Rates are 10–20\% lower than ChestX-ray. High color variation and artifacts (hair, rulers) reduce smoothness. Irregular lesion borders prove prone to certification failure. Certified pixels concentrate on lesion centers (pigmentation) rather than edges. Interestingly, certification inadvertently filters hair artifacts—a positive side effect where robustness guarantees improve real-world reliability by ignoring spurious features.

\begin{table}[h!]
\centering
\caption{IG certification rate at $\sigma=0.5$, $N=100$ (ResNet-18, Experiment 1)}
\begin{tabularx}{\linewidth}{|X|X|}
\hline
	extbf{Dataset} & \textbf{Certification rate} \\
\hline
ChestX-ray14 & 68\% \\
\hline
Brain MRI & 62\% \\
\hline
APTOS DR & 58\% \\
\hline
ISIC Skin & 52\% \\
\hline
\end{tabularx}
\label{tab:ig_cert_summary}
\end{table}

\subsection{Sample Size ($N$) Sensitivity Analysis}

Certification bounds tighten with increased Monte Carlo samples $N$. At $\sigma=0.5$ for Integrated Gradients, ChestX-ray achieves 54\% certification at $N=5$ versus 68\% at $N=100$—diminishing returns beyond $N=50$. At $N=5$, rates reach 75\% of $N=100$ values with 20× computation savings, suggesting $N=5$ provides practical default. Brain MRI shows similar pattern (49\% to 62\%), APTOS (44\% to 58\%), ISIC (39\% to 52\%). The recommendation balances coverage and efficiency at $N=50$.

\subsection{Certified Region Visualization and Analysis}

Certified attributions exhibit distinct spatial patterns across modalities. ChestX-ray certified regions focus on lung fields and cardiac silhouette, respecting anatomical boundaries. Brain MRI certifies tumor cores reliably while surrounding edema frequently fails certification. Fundus imaging shows major hemorrhages well-certified; microaneurysms frequently uncertified. ISIC lesion centers achieve certification while borders and artifacts remain uncertified. These patterns reflect the fundamental trade-off: certification requires smoothing that preserves global structures but obscures fine details critical for granular diagnosis.

\subsection{Computational Cost Analysis}

Gradient methods remain clinically feasible: Integrated Gradients (4.2 sec/image), GradCAM (3.1 sec/image), LRP (6.8 sec/image). Certifying 100 images requires 7–12 minutes per method. Perturbation methods become prohibitive: RISE consumes 52 sec/image (87 minutes for 100 images), Occlusion 98–101 sec/image. Total Experiment 1 computation across all methods, datasets, and parameter combinations consumed approximately 450 GPU-hours on NVIDIA A100 clusters. For clinical deployment, gradient methods enable offline analysis (e.g., certify diagnostic cases overnight); perturbation methods remain infeasible for real-time or near-real-time workflows.

\newpage

\section{Experiment 2: Cross-Architecture Comparison}

Experiment 2 evaluates all six architectures on a smaller subset (5 images per dataset, 20 total) using standard settings ($\sigma=0.5$, $N=50$, Integrated Gradients) to understand architecture effects on certification. ResNet-50 achieves highest average certification (61\%), followed by ResNet-34 (60\%) and ResNet-18 (58\%), with both EfficientNet variants and DenseNet at 56–58\%. Interestingly, EfficientNet-B1—achieving highest classification accuracy (89.7\% average)—paradoxically shows lowest certification (54\%), suggesting that higher-capacity models may overfit to spurious features harder to certify. Differences remain modest (±5–7\%), confirming that attribution method choice dominates architecture selection for certification. All architectures show consistent modality ranking: ChestX-ray > Brain MRI > APTOS > ISIC.

\begin{table}[h!]
\centering
\caption{Average certification rate by architecture (IG, $\sigma=0.5$, $N=50$, 20 images total)}
\begin{tabularx}{\linewidth}{|X|X|}
\hline
	extbf{Architecture} & \textbf{Avg certification} \\
\hline
ResNet-18 & 58\% \\
\hline
ResNet-34 & 60\% \\
\hline
ResNet-50 & 61\% \\
\hline
EfficientNet-B0 & 56\% \\
\hline
EfficientNet-B1 & 54\% \\
\hline
DenseNet-121 & 58\% \\
\hline
\end{tabularx}
\label{tab:arch_cert_summary}
\end{table}

\newpage

\section{Clinical Ground Truth Alignment Analysis}

Certified regions align reasonably well with clinical ground truth. For 30 brain MRI images with radiologist tumor annotations, certified Integrated Gradients attributions achieve 68\% Jaccard IoU (versus 72\% for uncertified), 81\% Dice coefficient, 76\% sensitivity (true positive rate), and 91\% specificity. The modest 3–4\% IoU loss reflects certification's smoothing effect. High specificity (91\%) suggests certified regions remain specific to actual tumors.

For ISIC skin lesion subset with segmentation masks, certified attributions show 62\% lesion localization IoU (versus 65\% uncertified), 71\% true positive rate, and 12\% false positive rate (versus 8\% uncertified). The slightly elevated false positive rate suggests certified attributions sometimes highlight background near lesion boundaries, but lesion-specific regions preserve clinical focus. Notably, certification filters spurious features: hair, rulers, and skin marking artifacts receive lower certified attribution—a beneficial side effect where robustness constraints inadvertently improve real-world performance.

APTOS fundus images show mixed results: certified attributions successfully localize 72\% of moderate/large hemorrhages, 45\% of microaneurysms, 38\% of hard exudates (texture-based, challenged by smoothing), and 89\% of the optic disc (large structure, naturally smooth). This reveals the critical trade-off: certification succeeds on coarse, clinically obvious structures but struggles with fine vasculature.

\newpage

\section{Grid Visualization and Certification (ISIC Dataset)}

To understand spatial focus and certification patterns in complex ISIC dermoscopy images, we generated 3×3 grid overlays computing attribution scores per cell (75×75 pixel regions). This reveals which image regions the model focuses on and how certification affects spatial attention.

Certified attributions show pronounced center focus compared to uncertified. Center cells achieve 68\% certification rates, adjacent cells 52\%, and corners 38\%. This pattern reflects lesion centering in dermoscopy protocols and smoothing's preservation of global structure while obscuring boundaries. Certified attributions increase center focus (+6% mean attribution) and significantly reduce edge/corner attention (21–22\% variance reduction). Grid-level localization shows certified attributions concentrate 68\% of importance on centers (versus 58\% uncertified), with slight lesion containment decrease (71\% versus 73\%) and notably lower boundary clarity (0.42 versus 0.58). The trade-off is acceptable for deployment: robustness and artifact filtering (hair, rulers) outweigh modest boundary precision loss. Sample grid visualizations show melanoma with uncertified diffuse heat versus certified concentration on pigmentation; benign nevus with uniform patterns; hair artifacts filtered in certified versions; and edge artifacts re-centered to lesion regions.

\newpage

\section{Summary of Key Results}

Certification effectiveness varies significantly by method and modality. Integrated Gradients achieves 52–68\% certification (highest), GradCAM 38–55\%, LRP 32–52\%, RISE 20–28\%, and Occlusion 12–22\%. Modality impact shows systematic pattern: grayscale ChestX-ray and Brain MRI outperform RGB APTOS and ISIC by ~15\%, with color channels and fine structures reducing certification coverage. Architecture impact remains modest (±5–7\%): ResNet-50 performs best, but attribution method choice dominates. Paradoxically, EfficientNet-B1 achieves highest accuracy (89.7\%) yet lowest certification (54\%), suggesting overfitting to spurious features.

Clinical alignment remains acceptable: certified regions achieve 62–71\% IoU with radiologist annotations, preserving major diagnostic structures while missing fine details. The trade-off is clear: smoothing required for robustness guarantees sacrifices boundary precision. Computational cost analysis confirms gradient methods (3–8 sec/image) clinically feasible for offline analysis; perturbation methods (45–120 sec) remain prohibitive for deployment. Total Experiment 1 consumed ~450 GPU-hours.

Artifact filtering emerges as unexpected benefit: certification inadvertently filters hair, rulers, and image padding artifacts, potentially improving real-world robustness beyond formal $\ell_2$ guarantees. This suggests certified XAI methods may prove more reliable in clinical practice than theoretical guarantees alone would predict.

\appendix

\section{Additional Visualizations}

[Placeholder for high-resolution certification heatmaps, grid overlays, and method comparison figures]

\end{document}
